La optimización de portafolios es un problema central en finanzas cuantitativas. Métodos clásicos como el enfoque de Markowitz proporcionan soluciones ampliamente aceptadas, mientras que alternativas más recientes como la Jerarquía de Riesgos (HRP) introducen enfoques basados en clustering y distancias de correlación.

La computación cuántica abre nuevas posibilidades para la representación y procesamiento de información financiera. En este contexto, este trabajo propone la implementación y validación experimental del algoritmo Quantum Hierarchical Risk Parity (QHRP), que integra codificación cuántica mediante un ansatz variacional con técnicas de ordenamiento jerárquico de riesgos. El objetivo es analizar si esta formulación permite obtener representaciones alternativas del espacio de correlaciones y comparar su comportamiento con el método clásico HRP.

Se presenta la formulación del ansatz cuántico, el esquema de codificación angular de datos financieros, la arquitectura del circuito y el flujo de procesamiento híbrido. Los experimentos se realizan en simuladores cuánticos y se comparan con implementaciones clásicas, evaluando la estabilidad de la codificación, la reproducibilidad de los resultados y las limitaciones del enfoque.

\palabrasclave{Computación cuántica, optimización de portafolios, Hierarchical Risk Parity, codificación cuántica, correlación financiera, validación experimental}